\section{A certified logarithmic invariant and exact cubic weights}
\label{sec:cubic-logarithmic}

This section retains $K=\mathbb Q(\omega)$,
$E=K(\alpha)$ with $\alpha^3=269$,
$\mathcal O=\mathbb Z_3[\omega]$ and $\lambda=\omega-1$.
The degree-zero logarithmic class group at three is denoted by
$\widetilde{\operatorname{Cl}}_3(E)$.  Degree zero is part of this
notation: the full logarithmic divisor-class group has an additional
free degree direction.  Write
\[
 W_e(E)=\ker\left(H^2(G_{E,S_3},\mathbb Z_3(e+1))
       \longrightarrow\bigoplus_{v\mid3}
                   H^2(E_v,\mathbb Z_3(e+1))\right) \quad(e\geq1).
\]
By AC-K10 and AC-K12 this is the wild \'etale kernel
$S_3(E,e)^\vee$, and
\[
 W_e(E)\simeq\mathcal O/\lambda^{a_e},\qquad
 a_e=\kappa_e-2(1+v_3(e))=m_e-1.                    \tag{KL.1}
\]
We also use $W_0(E)=\widetilde{\operatorname{Cl}}_3(E)$, the
weight-zero convention in the Iwasawa description below.

\begin{proposition}[AC-K13: the residual logarithmic invariant]
\label{AC-K13}
For every $e\geq1$,
\[
 \min\{2,m_e-1\}
 =\dim_{\mathbb F_3}\bigl(
       \widetilde{\operatorname{Cl}}_3(E)/3\bigr).       \tag{KL.2}
\]
Thus a residual logarithmic dimension zero would give $m_e=1$
for every $e$, dimension one would give $m_e=2$ for every $e$,
and dimension two gives $m_e\geq3$ for every $e$.
For the present field, the certified calculation described below gives
\[
 \operatorname{Cl}(E)\simeq\mathbb Z/7\oplus\mathbb Z/7,
 \qquad
 \widetilde{\operatorname{Cl}}_3(E)
       \simeq\mathbb Z/9\oplus\mathbb Z/3.             \tag{KL.3}
\]
Consequently the third alternative applies.  Neither theorem A
nor a conjecture about logarithmic class groups is used.
The group isomorphisms in (KL.3) are computer-assisted arithmetic
statements, with an unconditional certification command and its
returned result retained in the evidence files.
\end{proposition}
\begin{proof}
The precise known result used is J.-F.~Jaulent and A.~Michel,
\emph{Approche logarithmique des noyaux \'etales sauvages des corps
 de nombres}, J.~Number Theory \textbf{120} (2006), 72--91,
Theorem~3 and its complete proof on p.~7 of the author's 15-page
version:
\url{https://www.math.u-bordeaux.fr/~jjaulent/Articles/J44.pdf}.
Apply it with the auxiliary prime-to-three extension equal to the
identity extension of $E$ and with $r=1$.  Its cyclotomic
procyclicity hypothesis holds for the odd prime three, and
$\mu_3\subset E$.  It gives
\[
 W_e(E)/3\simeq\mu_3^{\otimes e}\otimes_{\mathbb Z_3}
                   \widetilde{\operatorname{Cl}}_3(E). \tag{KL.4}
\]
In particular their dimensions over $\mathbb F_3$ agree.
Since $3$ is a unit times $\lambda^2$ and the residue field of
$\mathcal O$ has three elements, (KL.1) gives
$\dim_{\mathbb F_3}(W_e/3)=\min(2,a_e)$.  AC-G26 now proves
(KL.2) and its three alternatives.  The source explicitly does
not assume Gross--Kuz'min, Schneider or Leopoldt conjectures
for its quotient isomorphism; no finiteness of a weight-zero
class group was inserted into this deduction.

Here are the exact input and certification for (KL.3).  With
$z=\omega-\alpha$, the resultant of
$Y^3-269$ and $(z+Y)^2+(z+Y)+1$ is
\[
 P(z)=z^6+3z^5+6z^4+545z^3+813z^2-804z+71824.
                                                               \tag{KL.5}
\]
The element $z$ generates $E$: indeed
\[
 \alpha=\frac{269-(2z+1)(z^2+z+1)}{3z(z+1)},
 \qquad \omega=z+\alpha,
\]
and $z(z+1)\ne0$.  Thus there is no unverified field
identification between the computational polynomial and $E$.
PARI/GP~2.19.1, built from the official source release, was run
on this exact polynomial with \texttt{bnfinit(P,1)}.
It returned discriminant $-141375086667$, class group
\texttt{[7,7]}, and class number $49$.
The complete command \texttt{bnfcertify(b)} returned \texttt{1};
no partial-certification flag was used.  Then
\texttt{bnflog(b,3)} returned
\texttt{[[9,3],[9,3],[]]}.
The first component is precisely the degree-zero logarithmic
class group.  The other components are its subgroup generated
at three and the ordinary $S_3$-class group, respectively.
The meaning and proof of this algorithm are K.~Belabas and
J.-F.~Jaulent, \emph{The logarithmic class group package in
PARI/GP}, Publ.~Math.~Besan\c con (2016), 5--18,
author version \S\S3.2--3.4, pp.~9--11, and \S4.1, p.~12:
\url{https://www.math.u-bordeaux.fr/~kbelabas/research/bnflog-2016.pdf}.
Its finite stopping criterion is proved in Lemma~3.8; the
initial class-group/unit data were unconditionally certified
here.  The reproducible script, complete raw output and binary
arithmetic data are
\texttt{evidence/cubic\_logarithmic/certify\_logclass.gp},
\texttt{certify\_logclass.out}, and
\texttt{cubic\_bnf.bin}.  The runtime build receipt is
\texttt{evidence/pari\_build.json}.  This proves the asserted
computer-assisted values and completes the residual conclusion.
\end{proof}

\begin{theorem}[AC-K14: finite-twist congruence and an exact infinite family]
\label{AC-K14}
Let $e,e'\geq0$ and $n\geq1$.  If
$e\equiv e'\pmod{3^{n-1}}$, then
\[
                   W_e(E)/3^n\simeq W_{e'}(E)/3^n.    \tag{KL.6}
\]
In particular, for every positive integer $e$ divisible by three,
\[
 \boxed{m_e=4,\qquad a_e=3,\qquad
        \kappa_e=2(1+v_3(e))+3.}                     \tag{KL.7}
\]
More generally, an exact value $a_e=a$ at one positive weight
fixes $a_{e'}=a$ and $m_{e'}=a+1$ at every positive weight
\[
                    e'\equiv e\pmod{3^{\lfloor a/2\rfloor}}.
                                                               \tag{KL.8}
\]
Thus the remaining weights can be investigated by finite
certificates on residue classes.  This theorem alone does not
give an upper bound on the two nonzero classes modulo three.
\end{theorem}
\begin{proof}
We spell out why (KL.6) uses a weaker condition than requiring
$\mu_{3^n}\subset E$.  Let $E_\infty/E$ be the cyclotomic
$\mathbb Z_3$-extension, $\Gamma_\infty=\langle\gamma\rangle$,
and let $X'$ be the Galois group of the maximal abelian pro-three
extension of $E_\infty$ unramified everywhere and split at
its primes above three.  The Schneider description used and
stated in the proof of Jaulent--Michel Theorem~3 is
\[
 W_e(E)=\bigl(\mathbb Z_3(e)\otimes_{\mathbb Z_3}X'\bigr)
                   _{\Gamma_\infty},
 \qquad
 W_0(E)=X'_{\Gamma_\infty}
             =\widetilde{\operatorname{Cl}}_3(E).    \tag{KL.9}
\]
The quotients are compact coinvariants.  At positive weights
this is the usual wild \'etale kernel occurring in (KL.1).
These known identifications, not a conjectural characteristic
ideal formula, are the only Iwasawa-theoretic input here.

Choose a basis of the Tate module and put
$u=\varepsilon_3(\gamma)\in1+3\mathbb Z_3$.
The quotient in (KL.9) modulo $3^n$ is the quotient of
$X'/3^nX'$ by the image of $u^e\gamma-1$.
Lifting the exponent in $1+3\mathbb Z_3$ shows
\[
 e\equiv e'\pmod{3^{n-1}}
       \quad\Longrightarrow\quad
                    u^e\equiv u^{e'}\pmod{3^n}.
\]
The two endomorphisms on $X'/3^nX'$ therefore agree, proving
(KL.6).  Images and quotients can equally be taken with closures:
they are continuous images of compact groups.  Changing the
chosen Tate basis only changes the displayed identification by
a unit.  If desired the isomorphism is compatible with the
cubic action by choosing its lift to fix all cyclotomic roots,
which is possible since $E\cap\mathbb Q(\mu_{3^\infty})=K$.

For $3\mid e$, take $e'=0$ and $n=2$.
By (KL.3), the right side of (KL.6) has order $27$.
By (KL.1), the left side has order
$3^{\min(a_e,4)}$.  Hence
$\min(a_e,4)=3$, which forces $a_e=3$.
Equation (KL.1) gives all of (KL.7).

Finally suppose $e,e'>0$ and $a_e=a$.
Take $n=\lfloor a/2\rfloor+1$, so that $2n>a$.
Under (KL.8), (KL.6) and the cyclic structure (KL.1) imply
\[
            \min(a_{e'},2n)=\min(a,2n)=a<2n.
\]
Consequently $a_{e'}=a$, and AC-G26 gives the equality of $m$.
This proves the general congruence assertion.
\end{proof}

\begin{remark}[The exact small computational gate]
The argument does not identify logarithmic-class vanishing
with ordinary three-rationality.  In fact the ordinary class
group has no three-part while the logarithmic class group in
(KL.3) is nonzero.  The three local places over three account
for two residual Brauer directions in AC-K12.
For the remaining weights, certifying $a_1$ and $a_2$ at values
less than four would settle both nonzero classes modulo three
by (KL.8).  A larger value requires finer residue classes.
Neither an uncertified class-group calculation nor a claim
about a regulator or an Iwasawa main conjecture is used as
such a certificate.
\end{remark}
